\documentclass{article}
\usepackage[UTF8]{ctex}
\usepackage{amsmath,amssymb,amsfonts}
\usepackage{enumitem}
\usepackage{geometry}
\geometry{a4paper, margin=1in}
% === 错题本排版优化（不参与题目标识） ===
\setlength{\parindent}{0pt}
\setlength{\parskip}{0.35em plus 0.10em minus 0.08em}
\setlength{\emergencystretch}{3em}
\setlist{itemsep=0.25em, topsep=0.40em, parsep=0pt}
\allowdisplaybreaks[2]
% === 错题本排版优化结束 ===
\begin{document}
\section*{合肥工业大学2024年数学分析真题}
\section*{107．合肥工业大学2024年数学分析试卷}

\subsection{题目1}
1．（10 分）求极限

\[
\lim _{n \rightarrow \infty} n^{2}\left(3^{\frac{1}{n}}-3^{\frac{1}{n+1}}\right)
\]

\subsection{题目2}
2．（10 分）求极限

\[
\lim _{x \rightarrow 0} \frac{1}{x^{4}}\left[\ln \left(1+\sin ^{2} x\right)-6(\sqrt[3]{2-\cos x}-1)\right]
\]

\subsection{题目3}
3．（10 分）判断数项级数 \(\sum_{n=2}^{\infty}(-1)^{n} \ln \left(1+\frac{1}{n}\right)\) 是否收敛，如果收敛，则判断是条件收敛还是绝对收敛．

\subsection{题目4}
4．（15 分）求极限

\[
\lim _{n \rightarrow \infty} \frac{\ln n}{\ln \left(1^{2023}+2^{2023}+\cdots+n^{2023}\right)}
\]

\subsection{题目5}
5．（15 分）设 \(f(x)\) 在 \([a, b]\) 上可积，且在 \(x=b\) 处连续，证明：

\[
\lim _{n \rightarrow \infty} \frac{n+1}{(b-a)^{n+1}} \int_{a}^{b}(x-a)^{n} f(x) \mathrm{d} x=f(b)
\]

\subsection{题目6}
6．（15 分）设函数 \(f(x)\) 在 \((-\infty,+\infty)\) 上有界且 \(f^{\prime \prime}(x) \geq 0\) ．证明：\(f(x)\) 为常值函数．

\subsection{题目7}
7．（15 分）求幂级数 \(\sum_{n=0}^{\infty}(n+2)(n+1) x^{n}\) 的收敛域及和函数．

\subsection{题目8}
8．（15 分）证明：\(F(x)=\int_{1}^{+\infty} \frac{y \sin x y}{1+y^{2}} \mathrm{~d} y\) 在 \((0,+\infty)\) 上连续．

\subsection{题目9}
9．（15 分）设有方程 \(\cos (x+y)+e^{y+x^{2}}-x^{3} y^{3}=2\) ．

\begin{itemize}
  \item[（1）] 证明：在点 \((0,0)\) 的某邻域内能确定唯一的隐函数 \(y=f(x)\) ．
  \item[（2）] 判断在 \((0,0)\) 的某邻域内能否确定隐函数 \(x=f(y)\) ？并说明理由．
\end{itemize}

\subsection{题目10}
10．（15 分）已知

\[
f(x, y)= \begin{cases}x y \frac{x^{2}-y^{2}}{x^{2}+y^{2}}, & (x, y) \neq(0,0) \\ 0, & (x, y)=(0,0)\end{cases}
\]

证明：\(f(x, y)\) 在 \((0,0)\) 处连续，并计算 \(f_{x}(0,0)\) 和 \(f_{x y}(0,0)\) ．

\subsection{题目11}
11．（15 分）求曲面积分

\[
\iint_{S} x \mathrm{~d} y \mathrm{~d} z+y \mathrm{~d} z \mathrm{~d} x+\left[(z+1)^{2}+e^{-\left(x^{2}+y^{2}\right)}\right] \mathrm{d} x \mathrm{~d} y
\]

其中 \(S\) 为 \(z=1-x^{2}-y^{2}(0 \leq z \leq 1)\) ，方向取下侧。
\end{document}
